\subsection{Matrix Equivalence and Similarity}
\begin{definition} \label{def:matrix-equivalence} \hlt{Matrix Equivalence}\\
For $A,B\in\operatorname{Mat}(m,n,\mathbb{F})$, write $A\sim_{\mathrm{eq}} B$ if $\exists$ $P\in\operatorname{Mat}(m,m,\mathbb{F})^{\times}$ and $Q\in\operatorname{Mat}(n,n,\mathbb{F})^{\times}$ s.t. $B=PAQ$.
\end{definition}

\begin{proposition} \label{prop:matrix-equivalence-basic-properties} \hlt{Matrix Equivalence and Elementary Operations}\\
The relation $\sim_{\mathrm{eq}}$ is an equivalence relation on $\operatorname{Mat}(m,n,\mathbb{F})$. Moreover, $A\sim_{\mathrm{eq}}B$ if and only if $B$ can be obtained from $A$ by finitely many elementary row and column operations.
\begin{proof}
Reflexivity follows from $A=I_mAI_n$. If $B=PAQ$ with $P,Q$ invertible, then $A=P^{-1}BQ^{-1}$, so symmetry holds. If $B=P_1AQ_1$ and $C=P_2BQ_2$, then $C=(P_2P_1)A(Q_1Q_2)$, and products of invertible matrices are invertible by Corollary \ref{cor:matrix-inverse-product-and-transpose}; hence transitivity holds.\\
Each elementary row operation is left multiplication by an invertible elementary matrix by Proposition \ref{prop:elementary-matrices-perform-row-operations}. Each elementary column operation is right multiplication by the transpose of an invertible elementary matrix, again by Proposition \ref{prop:elementary-matrices-perform-row-operations} and Corollary \ref{cor:matrix-inverse-product-and-transpose}. Thus any finite sequence of elementary row and column operations has the form $B=PAQ$ with $P,Q$ invertible, so $A\sim_{\mathrm{eq}}B$.\\
Conversely, suppose $A\sim_{\mathrm{eq}}B$, so $B=PAQ$ with $P,Q$ invertible. Lemma \ref{lem:invertible-matrices-row-reduce-to-identity} gives products of elementary matrices expressing $P$ and $Q$ as products of elementary factors. Left multiplication by the factors of $P$ is a sequence of elementary row operations, and right multiplication by the factors of $Q$ is a sequence of elementary column operations. Hence $B$ can be obtained from $A$ by finitely many elementary row and column operations.
\end{proof}
\end{proposition}

\begin{theorem} \label{thm:matrix-equivalence-same-linear-map} \hlt{Equivalent Matrices Represent the Same Linear Map}\\
Let $\dim\mathcal{V}=n$, $\dim\mathcal{W}=m$, and $A,B\in\operatorname{Mat}(m,n,\mathbb{F})$.\\
Then $A\sim_{\mathrm{eq}}B$ $\Longleftrightarrow$ $\exists\,T\in\LL(\mathcal{V},\mathcal{W})$ with $A=[T]_{\gamma\leftarrow\beta},\ B=[T]_{\gamma'\leftarrow\beta'}$ for some bases $\beta,\beta'$ of $\mathcal{V}$ and $\gamma,\gamma'$ of $\mathcal{W}$.
\begin{proof}
\begin{enumerate}[label=\roman*.]
\setlength{\itemsep}{0pt}
\item[$(\Rightarrow)$] Let $B=PAQ$ with $P,Q$ invertible. Choose bases $\beta,\gamma$ and define $T\in\LL(\mathcal{V},\mathcal{W})$ by $[T]_{\gamma\leftarrow\beta}=A$.\\
By Theorem \ref{thm:change-of-basis-matrices}, choose basis $\gamma'$ with $P=C_{\gamma'\leftarrow\gamma}$, and basis $\beta'$ with $Q^{-1}=C_{\beta'\leftarrow\beta}$. Then $Q=C_{\beta\leftarrow\beta'}$.\\
Then Theorem \ref{thm:two-sided-change-of-basis-for-linear-maps} gives $[T]_{\gamma'\leftarrow\beta'}=PAQ=B$.
\item[$(\Leftarrow)$] If $A=[T]_{\gamma\leftarrow\beta}$ and $B=[T]_{\gamma'\leftarrow\beta'}$, then $B=C_{\gamma'\leftarrow\gamma}AC_{\beta\leftarrow\beta'}$, so $A\sim_{\mathrm{eq}}B$.
\end{enumerate}
\hlt{Structural Commentary}: Equivalence is two-sided basis change.\\
\hlt{Key Insight}: Same map; different coordinates.
\end{proof}
\end{theorem}

\begin{theorem} \label{thm:rank-classifies-matrix-equivalence} \hlt{Rank Classifies Matrix Equivalence}\\
For $0\leq r\leq\min(m,n)$, set $J_r^{m,n}\defeq
\begin{bmatrix}
I_r & 0\\
0 & 0
\end{bmatrix}
\in\operatorname{Mat}(m,n,\mathbb{F})$.\\
For $A,B\in\operatorname{Mat}(m,n,\mathbb{F})$, $A\sim_{\mathrm{eq}}J_{\operatorname{rank}(A)}^{m,n}$ and $A\sim_{\mathrm{eq}}B\Longleftrightarrow \operatorname{rank}(A)=\operatorname{rank}(B)$.
\begin{proof}
\begin{enumerate}[label=\roman*.]
\setlength{\itemsep}{0pt}
\item[$(\Rightarrow)$] If $B=PAQ$ with $P,Q$ invertible, then $\operatorname{rank}(B)=\operatorname{rank}(A)$ by Theorem \ref{thm:full-rank-products-preserve-rank}.
\item[$(\Leftarrow)$] For any $X$ with $\operatorname{rank}(X)=r$, let $T_X:\mathbb{F}^n\to\mathbb{F}^m$ be $T_X(\mathbf{x})=X\mathbf{x}$.\\
Choose $\mathcal{M}_X$ with $\mathbb{F}^n=\ker(T_X)\oplus\mathcal{M}_X$; then $T_X|_{\mathcal{M}_X}:\mathcal{M}_X\to\operatorname{range}(T_X)$ is an isomorphism.\\
Choose a basis $(\mathbf{u}_1,\ldots,\mathbf{u}_r)$ of $\mathcal{M}_X$, extend it by a basis of $\ker(T_X)$ to $\beta_X$, and choose $\gamma_X$ beginning with $T_X\mathbf{u}_1,\ldots,T_X\mathbf{u}_r$. Then $[T_X]_{\gamma_X\leftarrow\beta_X}=J_r^{m,n}$, so $X\sim_{\mathrm{eq}}J_r^{m,n}$.\\
If $\operatorname{rank}(A)=\operatorname{rank}(B)=r$, then $A\sim_{\mathrm{eq}}J_r^{m,n}$ and $B\sim_{\mathrm{eq}}J_r^{m,n}$. By Proposition \ref{prop:matrix-equivalence-basic-properties}, symmetry gives $J_r^{m,n}\sim_{\mathrm{eq}}B$, and transitivity gives $A\sim_{\mathrm{eq}}B$.
\end{enumerate}
\hlt{Structural Commentary}: Rank is the invariant of two-sided equivalence.\\
\hlt{Key Insight}: Adapt bases to kernel and range.
\end{proof}
\end{theorem}

\begin{corollary} \label{cor:linear-map-rank-normal-form} \hlt{Rank Normal Form for Linear Maps}\\
Let $T\in\LL(\mathcal{V},\mathcal{W})$, $\dim\mathcal{V}=n$, $\dim\mathcal{W}=m$, and $\operatorname{rank}(T)=r$. Then $\exists$ bases $\beta,\gamma$ with $[T]_{\gamma\leftarrow\beta}=J_r^{m,n}$.
\begin{proof}
For initial bases $\beta_0,\gamma_0$, let $A=[T]_{\gamma_0\leftarrow\beta_0}$.\\
By Theorem \ref{thm:rank-classifies-matrix-equivalence}, $A\sim_{\mathrm{eq}}J_r^{m,n}$; Theorem \ref{thm:matrix-equivalence-same-linear-map} gives bases $\beta,\gamma$ with $[T]_{\gamma\leftarrow\beta}=J_r^{m,n}$.\\
\hlt{Structural Commentary}: Rank gives the two-basis normal form.\\
\hlt{Key Insight}: Rank normalizes the map.
\end{proof}
\end{corollary}

\begin{corollary} \label{cor:linear-map-equivalence-by-rank} \hlt{Linear Maps Classified by Rank}\\
For $S,T\in\LL(\mathcal{V},\mathcal{W})$, write $S\sim_{\mathrm{bc}}T$ if $\exists$ bases $\beta_S,\beta_T$ of $\mathcal{V}$ and $\gamma_S,\gamma_T$ of $\mathcal{W}$ such that $[S]_{\gamma_S\leftarrow\beta_S}=[T]_{\gamma_T\leftarrow\beta_T}$.\\
Then $S\sim_{\mathrm{bc}}T$ $\Longleftrightarrow$ $\operatorname{rank}(S)=\operatorname{rank}(T)$
\begin{proof}
\begin{enumerate}[label=\roman*.]
\setlength{\itemsep}{0pt}
\item[$(\Rightarrow)$] If $[S]_{\gamma_S\leftarrow\beta_S}=[T]_{\gamma_T\leftarrow\beta_T}$, then the two maps have the same matrix rank, hence the same rank.
\item[$(\Leftarrow)$] If $\operatorname{rank}(S)=\operatorname{rank}(T)=r$, Corollary \ref{cor:linear-map-rank-normal-form} gives bases with $[S]_{\gamma_S\leftarrow\beta_S}=J_r^{m,n}=[T]_{\gamma_T\leftarrow\beta_T}$.
\end{enumerate}
\hlt{Structural Commentary}: Rank is the invariant of basis-change equivalence.\\
\hlt{Key Insight}: Same rank means same normal form.
\end{proof}
\end{corollary}

\begin{theorem} \label{thm:change-of-basis-formula} \hlt{Change of Basis Formula}\\
Let $\mathcal{V}$ be finite-dimensional, let $T \in \LL(\mathcal{V})$, and let $\beta$ be a basis of $\mathcal{V}$.
\begin{enumerate}[label=(\roman*)]
\setlength{\itemsep}{0pt}
\item If $\gamma$ is a basis of $\mathcal{V}$ and $S=C_{\gamma \leftarrow \beta}$, then $[T]_{\gamma \leftarrow \gamma} = S[T]_{\beta \leftarrow \beta}S^{-1}$.
\item If $S\in\operatorname{Mat}(n,n,\mathbb{F})$ is invertible, then there is a basis $\gamma$ of $\mathcal{V}$ such that $[T]_{\gamma \leftarrow \gamma} = S[T]_{\beta \leftarrow \beta}S^{-1}$.
\end{enumerate}
\begin{proof}
\begin{enumerate}[label=(\roman*)]
\setlength{\itemsep}{0pt}
\item Suppose $\gamma$ is a basis and $S=C_{\gamma \leftarrow \beta}$. By Theorem \ref{thm:change-of-basis-matrices}, $S^{-1}=C_{\beta \leftarrow \gamma}$.\\
For every $\mathbf{v}\in\mathcal{V}$, $[T\mathbf{v}]_{\gamma} =S[T\mathbf{v}]_{\beta} =S[T]_{\beta \leftarrow \beta}[\mathbf{v}]_{\beta} =S[T]_{\beta \leftarrow \beta}S^{-1}[\mathbf{v}]_{\gamma}$.\\
Since $[\cdot]_{\gamma}$ is onto $\mathbb{F}^n$, the representing matrices are equal.
\item If $S$ is invertible, Theorem \ref{thm:change-of-basis-matrices} supplies a basis $\gamma$ with $S=C_{\gamma \leftarrow \beta}$, and part (i) applies.
\end{enumerate}
\hlt{Structural Commentary}: Changing basis conjugates operator matrices.\\
\hlt{Key Insight}: Similarity is a coordinate rewrite.
\end{proof}
\end{theorem}

\begin{definition} \label{def:similar-matrices} \hlt{Similarity}\\
Matrices $A,B \in \operatorname{Mat}(n,n,\mathbb{F})$ are similar if there exists an invertible matrix $P$ such that $B = PAP^{-1}$.
\end{definition}

\begin{proposition} \label{prop:similarity-equivalence-relation} \hlt{Similarity Is an Equivalence Relation}\\
Similarity is an equivalence relation on $\operatorname{Mat}(n,n,\mathbb{F})$.
\begin{proof}
For any $A$, $A=I_nAI_n^{-1}$, so similarity is reflexive. If $B=PAP^{-1}$, then $A=P^{-1}BP$, so it is symmetric. If $B=PAP^{-1}$ and $C=QBQ^{-1}$, then $C=Q(PAP^{-1})Q^{-1}=(QP)A(QP)^{-1}$ so it is transitive.\\
\hlt{Structural Commentary}: Similarity groups matrices by change of basis.\\
\hlt{Key Insight}: Identity, inverse, and composition give the three equivalence properties.
\end{proof}
\end{proposition}

\begin{perspective} \hlt{Similarity as Same Operator, New Coordinates}\\
Similarity does not mean two operators are merely related; it means they are the same operator written in two different bases. This is the structural content behind change of basis.
\end{perspective}

\begin{perspective} \label{per:similarity-invariants-and-canonical-forms} \hlt{Similarity Invariants and Canonical Forms}\\
Similarity is equality up to change of basis. Trace, determinant, characteristic polynomials, and minimal polynomials are similarity invariants by Theorems \ref{thm:similarity-preserves-trace-and-determinant}, \ref{thm:similarity-preserves-characteristic-polynomial}, and \ref{thm:similarity-preserves-minimal-polynomial}, but they are not generally complete.\\
Canonical forms, such as the Jordan form of Theorem \ref{thm:complex-jordan-form}, choose one structured representative of each similarity class when possible.
\end{perspective}

\begin{remark} \label{rem:similarity-convention} \hlt{Similarity Convention}\\
The forms $PAP^{-1}$ and $M^{-1}AM$ are the same convention with $P=M^{-1}$. In Theorem \ref{thm:change-of-basis-formula}, this means choosing opposite coordinate-change directions.
\end{remark}

\begin{figure}[H]
  \centering
  \begin{tikzpicture}[
  >=Stealth,
  every node/.style={font=\small},
  coord/.style={draw=blue!60!black, fill=blue!7, rounded corners=2pt,
    minimum width=1.85cm, minimum height=0.64cm, align=center},
  coordb/.style={coord, fill=teal!7, draw=teal!60!black},
  map/.style={->, line width=0.78pt, blue!70!black},
  change/.style={->, line width=0.72pt, gray!70!black}
]
  \node[coord] (b0) at (0,0.88) {$[\mathbf{v}]_{\beta}$};
  \node[coord] (b1) at (4.05,0.88) {$[T\mathbf{v}]_{\beta}$};
  \node[coordb] (g0) at (0,-0.88) {$[\mathbf{v}]_{\gamma}$};
  \node[coordb] (g1) at (4.05,-0.88) {$[T\mathbf{v}]_{\gamma}$};

  \draw[map] (b0) -- node[above] {$A=[T]_{\beta\leftarrow\beta}$} (b1);
  \draw[map] (g0) -- node[below] {$B=[T]_{\gamma\leftarrow\gamma}$} (g1);
  \draw[change] (b0) -- node[left] {$P=C_{\gamma\leftarrow\beta}$} (g0);
  \draw[change] (b1) -- node[right] {$P$} (g1);
\end{tikzpicture}

  \caption{One operator in two bases: $PA=BP$, so $B=PAP^{-1}$.}
  \label{fig:similarity-change-of-basis-bridge}
\end{figure}

\begin{corollary} \label{cor:similar-iff-same-linear-operator} \hlt{Similarity as One Operator in Two Bases}\\
For $A,B\in\operatorname{Mat}(n,n,\mathbb{F})$, the following are equivalent:
\begin{enumerate}[label=(\roman*)]
\setlength{\itemsep}{0pt}
\item $A$ and $B$ are similar.
\item There are an $n$-dim vector space $\mathcal{V}$, bases $\beta,\gamma$ of $\mathcal{V}$, and $T\in\LL(\mathcal{V})$ such that $A=[T]_{\beta \leftarrow \beta}$ and $B=[T]_{\gamma \leftarrow \gamma}$.
\end{enumerate}
\begin{proof}
\begin{enumerate}[label=(\roman)*.]
\setlength{\itemsep}{0pt}
\item[$(i)$] $\Rightarrow (ii)$: Suppose $B=PAP^{-1}$ with $P$ invertible.\\
Let $\mathcal{V}=\mathbb{F}^n$, let $\beta$ be the standard basis, and let $T_A(\mathbf{x})=A\mathbf{x}$. By Theorem \ref{thm:change-of-basis-matrices}, choose $\gamma$ with $P=C_{\gamma \leftarrow \beta}$. Then $[T_A]_{\beta \leftarrow \beta}=A$, and Theorem \ref{thm:change-of-basis-formula} gives $[T_A]_{\gamma \leftarrow \gamma}=B$.
\item[$(ii)$] $\Rightarrow (i)$: If $A=[T]_{\beta \leftarrow \beta}$ and $B=[T]_{\gamma \leftarrow \gamma}$, then $B=C_{\gamma \leftarrow \beta}AC_{\beta \leftarrow \gamma}$ with $C_{\gamma \leftarrow \beta}$ invertible.
\end{enumerate}
\hlt{Structural Commentary}: Similar matrices are one operator in different bases.\\
\hlt{Key Insight}: Conjugation records a change of basis.
\end{proof}
\end{corollary}

\begin{definition} \label{def:similar-linear-operators} \hlt{Similarity of Operators}\\
For $T\in\LL(\mathcal{V})$ and $S\in\LL(\mathcal{W})$, $T\sim S$ $\Longleftrightarrow$ $\exists\,\Phi:\mathcal{V}\xrightarrow{\sim}\mathcal{W}$ such that $ S=\Phi T\Phi^{-1}$.\\
If $\mathcal{V}=\mathcal{W}$, then $\Phi\in\operatorname{Aut}(\mathcal{V})$ and similarity is conjugacy in $\LL(\mathcal{V})$.
\end{definition}

\begin{proposition} \label{prop:operator-similarity-and-common-matrices} \hlt{Operator Similarity and Common Matrix Representatives}\\
Let $\dim\mathcal{V}=\dim\mathcal{W}=n$, $T\in\LL(\mathcal{V})$, and $S\in\LL(\mathcal{W})$.\\
Then $T\sim S$ $\Longleftrightarrow$ $\exists\,A\in\operatorname{Mat}(n,n,\mathbb{F})$ with $A=[T]_{\beta\leftarrow\beta}=[S]_{\gamma\leftarrow\gamma}$ for some bases $\beta$ of $\mathcal{V}$ and $\gamma$ of $\mathcal{W}$.
\begin{proof}
\begin{enumerate}[label=\roman*.]
\setlength{\itemsep}{0pt}
\item[$(\Rightarrow)$] Let $S=\Phi T\Phi^{-1}$ and $\beta=(\mathbf{v}_1,\ldots,\mathbf{v}_n)$. Put $\gamma=(\Phi\mathbf{v}_1,\ldots,\Phi\mathbf{v}_n)$.\\
Then $\gamma$ is a basis and, for all $j$, $S(\Phi\mathbf{v}_j)=\Phi T\mathbf{v}_j$, so $[S]_{\gamma\leftarrow\gamma}=[T]_{\beta\leftarrow\beta}$.
\item[$(\Leftarrow)$] Suppose $[T]_{\beta\leftarrow\beta}=[S]_{\gamma\leftarrow\gamma}$. Define $\Phi:\mathcal{V}\to\mathcal{W}$ by $\Phi(\beta_j)=\gamma_j$.\\
Then $\Phi$ is an isomorphism and $\Phi T\beta_j=S\Phi\beta_j$ for every $j$, hence $\Phi T=S\Phi$. Thus $S=\Phi T\Phi^{-1}$.
\end{enumerate}
\hlt{Structural Commentary}: Similarity is equality of matrices after identifying bases.\\
\hlt{Key Insight}: A common matrix is an operator-level conjugacy written in bases.
\end{proof}
\end{proposition}

\begin{theorem} \label{thm:similarity-preserves-scalar-shifts} \hlt{Similarity and Scalar Shifts}\\
Let $A,B\in\operatorname{Mat}(n,n,\mathbb{F})$.
\begin{enumerate}[label=(\roman*)]
\setlength{\itemsep}{0pt}
\item If $A$ and $B$ are similar, then $A-\lambda I_n$ and $B-\lambda I_n$ are similar for every $\lambda\in\mathbb{F}$.
\item If $A-\lambda I_n$ and $B-\lambda I_n$ are similar for some $\lambda\in\mathbb{F}$, then $A$ and $B$ are similar.
\end{enumerate}
\begin{proof}
\begin{enumerate}[label=(\roman*)]
\setlength{\itemsep}{0pt}
\item If $B=PAP^{-1}$ with $P$ invertible, then $B-\lambda I_n=PAP^{-1}-\lambda PP^{-1}=P(A-\lambda I_n)P^{-1}$.
\item If $B-\lambda I_n=P(A-\lambda I_n)P^{-1}$ with $P$ invertible, then adding $\lambda I_n=\lambda PP^{-1}$ gives $B=PAP^{-1}$.
\end{enumerate}
\hlt{Structural Commentary}: Scalar shifts commute with change of basis.\\
\hlt{Key Insight}: The identity matrix is basis-independent.
\end{proof}
\end{theorem}
